\input{../tex-inputs/latex-korrekturansicht-vorspann}

\section[Berta Zuckerkandl an Arthur Schnitzler, {[}13. 6. 1911?{]}]{L03996 Berta Zuckerkandl an Arthur Schnitzler, {[}13. 6. 1911?{]}}
\nopagebreak\mylabel{L03996v}
\rehead{ }\normalsize\beginnumbering\briefempfaengerindex{Schnitzler, Arthur@\textsc{Schnitzler, Arthur}!zzzZuckerkandl, Berta@\emph{von Berta Zuckerkandl}!1911-06-131@{{[}13. 6. 1911?{]}}|(be}
\toendnotes[C]{\smallbreak\pagebreak[2]}
\correspDesc{Versand  durch Berta Zuckerkandl am [13. 6. 1911?] in Wien
\newline{}Erhalt  durch Arthur Schnitzler am [13. 6. 1911?] in Wien}\toendnotes[C]{\smallbreak}
\Standort{CUL, Schnitzler, B 200.}
\physDesc{Brief, 1 Blatt, 3 Seiten, 643 Zeichen
\newline{}Handschrift: schwarze Tinte, lateinische Kurrent
\newline{}Schnitzler: mit Bleistift beschriftet: »\textsc{Zucker}{[}kandl{]}« }\toendnotes[C]{\smallbreak}
\pstart{}{\pb}Hochverehrter Herr Doktor!\pend\vspace{0.5em}
\pstart
           Ich war sehr leidend und arbeitsunfähig. Auch hatte ich nichts Weiteres aus \textcolor{pink}{Paris}\oindex{Paris@\textbf{Paris}, \emph{Hauptstadt}|pw}{}\ledrightnote{\textcolor{pink}{Paris}} gehört. Nun erhalte ich eben einen \label{K_L03996-1v}\edtext{Brief meiner \textcolor{blue}{Schwester}\pwindex{Clemenceau, Sophie 25.\,5.\,1862 – 24.\,9.\,1937@\textsc{Clemenceau, Sophie} (25.\,5.\,1862 – 24.\,9.\,1937)|pwv}{}\ledrightnote{{$\rightarrow$}\emph{\textcolor{blue}{Sophie Clemenceau}}}}{\lemma{\textnormal{\emph{Brief meiner Schwester}}}\Cendnote{\textnormal{nicht überliefert}}}\label{K_L03996-1} den ich Ihnen
               vorlesen muss — da ich \label{K_L03996-2v}\edtext{Samstag}{\lemma{\textnormal{\emph{Samstag}}}\Cendnote{\textnormal{17. 6. 1911}}}\label{K_L03996-2}{ }Früh nach \textcolor{pink}{Paris}\oindex{Paris@\textbf{Paris}, \emph{Hauptstadt}|pw}{}\ledrightnote{\textcolor{pink}{Paris}} reise. {\pb}Und zwar wäre es \label{K_L03996-3v}\edtext{dringend dass ich Sie heute}{\lemma{\textnormal{\emph{dringend … heute}}}\Cendnote{\textnormal{Für den 13. 6. 1911 notiert \textcolor{blue}{Schnitzler} im \emph{\textcolor{green}{Tagebuch}\pwindex{Schnitzler, Arthur 15. 5. 1862 Wien – 21. 10. 1931 ebd.@\textsc{Schnitzler, Arthur} (15. 5. 1862 Wien – 21. 10. 1931 ebd.), \emph{Schriftsteller, Mediziner}!Tagebuch@\strich\emph{Tagebuch}|pwk}}
                  für den Nachmittag: »Frau Hofr. \textcolor{blue}{Zuckerkandl}\pwindex{Zuckerkandl, Berta 13.\,4.\,1864 Wien – 16.\,10.\,1945 Paris@\textsc{Zuckerkandl, Berta} (13.\,4.\,1864 Wien – 16.\,10.\,1945 Paris), \emph{Schriftstellerin, Journalistin, Übersetzerin}|pw}; in Sachen \textcolor{green}{Medardus}\pwindex{Schnitzler, Arthur 15. 5. 1862 Wien – 21. 10. 1931 ebd.@\textsc{Schnitzler, Arthur} (15. 5. 1862 Wien – 21. 10. 1931 ebd.), \emph{Schriftsteller, Mediziner}!junge Medardus. Dramatische Historie in einem Vorspiel und fünf Aufzügen@\strich\emph{Der junge Medardus. Dramatische Historie in einem Vorspiel und fünf Aufzügen}|pw}
                     für \textcolor{pink}{Paris}\oindex{Paris@\textbf{Paris}, \emph{Hauptstadt}|pw}.« Das entspricht \textcolor{blue}{Zuckerkandls}\pwindex{Zuckerkandl, Berta 13.\,4.\,1864 Wien – 16.\,10.\,1945 Paris@\textsc{Zuckerkandl, Berta} (13.\,4.\,1864 Wien – 16.\,10.\,1945 Paris), \emph{Schriftstellerin, Journalistin, Übersetzerin}|pwk} im vorliegenden Brief geäußertem
                  Wunsch, ihn noch am selben Tag zu sprechen zu wollen.}}}\label{K_L03996-3} spräche – denn ich
               soll ein \textcolor{pink}{französisches}\oindex{Frankreich@\textbf{Frankreich}|pw}{}\ledrightnote{\textcolor{pink}{Frankreich}}{ }\label{K_L03996-4v}\edtext{\textcolor{green}{Scenarium}\pwindex{Schnitzler, Arthur 15. 5. 1862 Wien – 21. 10. 1931 ebd.@\textsc{Schnitzler, Arthur} (15. 5. 1862 Wien – 21. 10. 1931 ebd.), \emph{Schriftsteller, Mediziner}!junge Medardus. (Zur Übersetzung ins Französische)@\strich\emph{Der junge Medardus. (Zur Übersetzung ins Französische)}|pwv}{}\ledrightnote{{$\rightarrow$}\emph{\textcolor{green}{Der junge Medardus. (Zur Übersetzung ins Französische)}}} des »\textcolor{green}{Medardus}\pwindex{Schnitzler, Arthur 15. 5. 1862 Wien – 21. 10. 1931 ebd.@\textsc{Schnitzler, Arthur} (15. 5. 1862 Wien – 21. 10. 1931 ebd.), \emph{Schriftsteller, Mediziner}!junge Medardus. Dramatische Historie in einem Vorspiel und fünf Aufzügen@\strich\emph{Der junge Medardus. Dramatische Historie in einem Vorspiel und fünf Aufzügen}|pw}{}\ledrightnote{\textcolor{green}{Der junge Medardus. Dramatische Historie in einem Vorspiel und fünf Aufzügen}}«}{\lemma{\textnormal{\emph{Scenarium des »Medardus«}}}\Cendnote{\textnormal{Vgl. Arthur Schnitzler an Berta Zuckerkandl, 20. 6. 1911.}}}\label{K_L03996-4} machen – eine grosse Arbeit – für
               welche ich Ihre Hilfe brauchte.\pend
           
\pstart
           – Wann könnte ich Sie nun heute länger
                  {\pb}sprechen. Ausser von halb
                  fünf – bis halb sechs Uhr wo ich bereits von meinem \textcolor{blue}{Arzt}\pwindex{?? [Arzt von Berta Zuckerkandl] @\textsc{?? [Arzt von Berta Zuckerkandl]}|pwv}{}\ledrightnote{{$\rightarrow$}\emph{\textcolor{blue}{?? [Arzt von Berta Zuckerkandl]}}} erwartet werde – bin
               ich ganz frei.\pend
           
\pstart
           Bitte gütigst \label{K_L03996-5v}\edtext{Antwort}{\lemma{\textnormal{\emph{Antwort}}}\Cendnote{\textnormal{nicht überliefert. Eventuell erfolgte diese telefonisch?}}}\label{K_L03996-5}. Geht es heute nicht – vielleicht dann
                  morgen{ }Vormittag?\pend
           \pstart Beste Grüsse \spacefill\mbox{B. Zuckerkandl}\pend{}\selectlanguage{ngerman}\endnumbering\briefempfaengerindex{Schnitzler, Arthur@\textsc{Schnitzler, Arthur}!zzzZuckerkandl, Berta@\emph{von Berta Zuckerkandl}!1911-06-131@{{[}13. 6. 1911?{]}}|)be}\mylabel{L03996h}  \newcommand{\dateiname}{L03996}\newcommand{\titel}{Berta Zuckerkandl an Arthur Schnitzler, [13. 6. 1911?]}\newcommand{\editorInnen}{Herausgegeben von Jahnke, SelmaMüller, Martin Anton}\input{../tex-inputs/latex-korrekturansicht-abspann}
